% Source: source/普通高考/2015/2015重庆文.pdf, page 1, question 8.
% pdfimages -list found no embedded bitmap; the program flowchart is vector artwork.
% Grid evidence: tmp/redraw_flowchart_2015_chongqing_liberal/grid/page1_grid100.jpg
% and tmp/redraw_flowchart_2015_chongqing_liberal/crop/q8_block_grid50.jpg.
% Node-edge list: start -> s=0,k=0 -> test(k<8); test(是) -> k=k+2
% -> s=s+1/k -> loop back to the central connector before the test;
% test(否) -> output s -> finish. No other connectors are present.
% All segments are solid directed connectors; node text is locally zero-indented
% and centered. The drawing is schematic and preserves the source topology.

\begin{tikzpicture}[
    >=Stealth,
    x=1cm,
    y=0.88cm,
    line width=0.45pt,
    line cap=round,
    line join=round,
    font=\normalsize,
    flowbox/.style={draw, anchor=center, minimum height=0.68cm,
        inner xsep=4pt, inner ysep=2.5pt, text centered,
        execute at begin node={\setlength{\parindent}{0pt}}},
    terminal/.style={flowbox, rounded corners=0.14cm, minimum width=1.25cm},
    process/.style={flowbox, minimum width=2.45cm},
    decision/.style={flowbox, diamond, aspect=3.2, minimum width=2.95cm,
        minimum height=1.00cm},
    io/.style={flowbox, trapezium, minimum width=1.60cm, minimum height=0.78cm,
        trapezium left angle=70, trapezium right angle=110},
]
    \path[use as bounding box] (-1.60,0.95) rectangle (3.95,11.10);

    \node[terminal] (start) at (0,10.65) {开始};
    \node[process, minimum width=2.65cm] (init) at (0,9.35) {$s=0,k=0$};
    \node[decision] (test) at (0,5.10) {$k<8$};
    \node[process, minimum width=1.95cm] (increment) at (2.55,6.05) {$k=k+2$};
    \node[process, minimum width=2.35cm, minimum height=0.86cm] (add)
        at (2.55,7.45) {$s=s+\dfrac{1}{k}$};
    \node[io] (output) at (0,3.20) {输出 $s$};
    \node[terminal] (finish) at (0,1.85) {结束};

    % Main downward path and the terminating branch.  The loop return first
    % meets this trunk, followed by one separate arrow into the test diamond.
    \draw[->] (start.south) -- (init.north);
    \coordinate (merge-level) at (0,8.95);
    \draw (init.south) -- (merge-level);
    \draw[->] (merge-level) -- (test.north);
    \draw[->] (test.south) -- (output.north);
    \draw[->] (output.south) -- (finish.north);

    % 是 branch: increment k, accumulate s, then return before the test.
    \draw[->] (test.east) -- (2.55,5.10) -- (2.55,5.57) -- (increment.south);
    \draw[->] (increment.north) -- (add.south);
    \coordinate (merge-pre) at (0.20,8.95);
    \draw[->] (add.north) -- (2.55,8.95) -- (merge-pre);
    \draw (merge-pre) -- (merge-level);
    \node[anchor=south west, text centered, inner xsep=1pt, inner ysep=1pt] at (0.92,5.23) {是};
    \node[anchor=west, text centered, inner xsep=1pt, inner ysep=1pt] at (0.14,4.42) {否};
\end{tikzpicture}
